% !TEX root = main.tex
%%%%%%%%%%%%%%%%%%%%%%%%%%%%%%%%%%%%%%%%%%%%%%%%%%%%%%%%%%%%%%%%%%%%%%%%%%%%%%%%
%  Beyond Scale: The Diversity Coefficient as a Data Quality Metric
%  TMLR submission (Transactions on Machine Learning Research)
%
%  Content is shared with paper_latex/DMLR_2026_BeyondScale/ (same modular
%  section files); this root retargets it to the TMLR style (tmlr.sty).
%
%  Style:   tmlr.sty  (double-blind submission mode; use [accepted] for camera-ready)
%  Bib:     dmlr_bib_refs.bib  (tmlr.bst)
%  Figures: plots/
%%%%%%%%%%%%%%%%%%%%%%%%%%%%%%%%%%%%%%%%%%%%%%%%%%%%%%%%%%%%%%%%%%%%%%%%%%%%%%%%

\documentclass[10pt]{article}
\usepackage{tmlr}
% For the camera-ready, instead use: \usepackage[accepted]{tmlr}
% For a de-anonymized preprint, instead use: \usepackage[preprint]{tmlr}

% Optional math commands from https://github.com/goodfeli/dlbook_notation.
\input{math_commands.tex}

\usepackage{hyperref}
\usepackage{url}
\usepackage{graphicx}
\usepackage{booktabs}
\usepackage{tabularx}
\usepackage{algorithm}
\usepackage{algorithmic}
% algorithmic.sty defines \AND (logical-and) at each \begin{algorithmic}; tmlr.sty
% defines \AND as the author separator (needed by \maketitle). They clash only in
% the body, so we undefine \AND right after \maketitle (below), leaving the author
% block intact and letting each algorithmic environment define its own \AND.
% \acks is a dmlr2e macro absent from tmlr.sty. Under double-blind TMLR
% submission, acknowledgments are omitted (they identify authors); restore for
% the camera-ready by redefining \acks to emit a \subsubsection*{Acknowledgments}.
\newcommand{\acks}[1]{}

\title{Beyond Scale: The Diversity Coefficient as a Data Quality Metric for Variability in Natural Language Data}

% Authors must not appear in the submitted version. They are hidden by tmlr.sty
% as long as the package is used without the [accepted] or [preprint] options.
\author{%
Brando Miranda\thanks{Correspondence to: \texttt{brando9@stanford.edu}\,\;and\;\texttt{sanmi@stanford.edu}}\space\space$^{1}$
\And
Alycia Lee$^{1}$
\And
Sudharsan Sundar$^{1}$
\AND
Elyas Obbad$^{1}$
\And
Allison Casasola$^{1}$
\And
Rylan Schaeffer$^{1}$
\AND
Sanmi Koyejo$^{1,*}$
\\[0.8em]
$^{1}$Department of Computer Science, Stanford University, Stanford, CA 94305, USA\\
\texttt{\{brando9, alylee15, sjsundar\}@stanford.edu}\\
\texttt{\{eobbad, allyc, rschaef, sanmi\}@stanford.edu}
}

\def\month{07}  % Insert correct month for camera-ready version
\def\year{2026} % Insert correct year for camera-ready version
\def\openreview{\url{https://openreview.net/forum?id=3ndFgcqAic}} % Insert correct link for camera-ready version

\begin{document}

\maketitle
% Release tmlr's author-separator \AND now that \maketitle is done, so the
% algorithmic environments in the body can define their own logical-and \AND.
\makeatletter\let\AND\@undefined\makeatother

\begin{abstract}
Current trends in pre-training Large Language Models (LLMs) primarily focus on the scaling of model and dataset size.
While the \textit{quality} of pre-training data is considered an important factor for training powerful LLMs, it remains a nebulous concept that has not been rigorously characterized.
To this end, we propose a formalization of one key aspect of data quality -- measuring the \textit{variability} of natural language data -- specifically via a measure we call the diversity coefficient.
Our empirical analysis shows that the proposed diversity coefficient aligns with the intuitive properties of diversity and variability,
e.g., it increases as the number of latent concepts increases.
Then, we measure the diversity coefficient of publicly available pre-training datasets and demonstrate that their formal diversity is high compared to theoretical lower and upper bounds.
Finally, we conduct a comprehensive set of controlled \textit{interventional} experiments with GPT-2 and LLaMAv2 that demonstrate the diversity coefficient of pre-training data characterizes useful aspects of downstream model evaluation performance---totaling 44 models of various sizes (51M to 7B parameters).
We conclude that our formal notion of diversity is an important aspect of data quality that captures variability and causally leads to improved evaluation performance.
\end{abstract}

%%%%%%%%%%%%%%%%%%%%%%%%%%%%%%%%%%%%%%%%%%%%%%%%%%%%%%%%%%%%%%%%%%%%%%%%%%%%%%%%
% Main sections (shared modular files)
%%%%%%%%%%%%%%%%%%%%%%%%%%%%%%%%%%%%%%%%%%%%%%%%%%%%%%%%%%%%%%%%%%%%%%%%%%%%%%%%
\input{01_introduction}

\input{02_method}

\input{03_experiments}

\input{04_related_work}

\input{05_discussion}

%%%%%%%%%%%%%%%%%%%%%%%%%%%%%%%%%%%%%%%%%%%%%%%%%%%%%%%%%%%%%%%%%%%%%%%%%%%%%%%%
\subsubsection*{Broader Impact Statement}
This work proposes the diversity coefficient as a quantitative metric for data quality in LLM pre-training.
By making data diversity formally measurable, we enable researchers and practitioners to design more principled and transparent data curation pipelines.
A potential concern is that diversity metrics could be gamed or misused to justify data collection practices that raise privacy or copyright issues; we encourage responsible use alongside appropriate data governance.
All datasets used in this study are publicly available.

%%%%%%%%%%%%%%%%%%%%%%%%%%%%%%%%%%%%%%%%%%%%%%%%%%%%%%%%%%%%%%%%%%%%%%%%%%%%%%%%
\input{98_acknowledgments}

%%%%%%%%%%%%%%%%%%%%%%%%%%%%%%%%%%%%%%%%%%%%%%%%%%%%%%%%%%%%%%%%%%%%%%%%%%%%%%%%
\bibliographystyle{tmlr}
\bibliography{dmlr_bib_refs}

%%%%%%%%%%%%%%%%%%%%%%%%%%%%%%%%%%%%%%%%%%%%%%%%%%%%%%%%%%%%%%%%%%%%%%%%%%%%%%%%
\appendix
\input{99_appendix}

\end{document}
